\section{Conclusion}
Fresh randomness is useful only to the extent that it creates fresh observable uncertainty. In \scheme{}, a legal repeated-token request makes the first-layer Value term invariant to row permutations. The remaining marker visits a small orbit whose rank-one differences identify an equivalent demixer. Across ten models and twelve configurations, the attack recovers 263,029 of 263,904 head--token opportunities, while wrong-secret controls remain between 0 and 0.625\%. The attack also consumes the official protected-cache output and transfers across requests within a secret epoch.

The result is strong but bounded: it recovers blockwise token multisets, begins with cache observation, and does not establish deployment prevalence or a complete repair. A limited closed-set study shows why those multisets can matter, while request-specific refresh blocks the tested cross-request transfer without proving general security. The broader design rule is to analyze the observable orbit induced jointly by randomness and reachable inputs. Counting random choices without that quotient can substantially overstate protection.
